\section{Related Work}
\label{sec:related_work}
\paragraph{Word Embeddings}
Word embeddings serve as a fundamental building block for a broad range of NLP applications~\cite{sentiment,machinetranslation,realtionextraction}
%, such as sentiment classification~\cite{sentiment}, relation extraction~\cite{realtionextraction}, and machine translation~\cite{machinetranslation}, etc.
 and various approaches~\cite{mikolov,glove,yoav} have been proposed for training the word vectors. % by leveraging the pair-wise co-occurrences of words observed in unannotated language corpus.
Improvements have been made by leveraging semantic lexicons and morphology~\cite{luong2013better,faruqui2014}, disambiguating multiple senses~\cite{vsuster2016bilingual,arora2018linear,upadhyay2017beyond}, and modeling contextualized information by deep neural networks~\cite{elmo}.  
% for word embedding models have been made by~\cite{faruqui2014,yoav,elmo}, 
%\newcite{faruqui2014} improved the semantic quality of word embeddings by adding relational information from semantic lexicons.
%\newcite{yoav} discussed how system design choice and parameter optimization affect the performance of word embeddings.
However, none of these works attempts to tackle the problem of stereotypes exhibited in embeddings.

% The first constructive implementation for word representations is proposed in~\cite{mikolov}. 
% Two basic models are described, Skip-gram and Continuous Bag-of-Words, to learn the representation of words.
%\zy{@TODO: I think it lacks a transition sentence}
%GloVe~\cite{glove} claims to enhance the word embedding by leveraging the pair-wise co-occurrence features of words together with a global log-bilinear regression model.
\paragraph{Stereotype Analysis}
Implicit stereotypes have been observed in applications such as online advertising systems~\cite{sweeney2013discrimination}, web search~\cite{kay2015unequal}, and online reviews~\cite{wallace2016jerk}.
Besides, \newcite{jieyu} and \newcite{Rachel18} show that  coreference resolution systems are gender biased. The systems can successfully predict the link between ``the president'' with male pronoun but fail with the female one.  \newcite{rudinger2017social} use point-wise mutual information
to test the SNLI~\cite{snli:emnlp2015} corpus and demonstrate gender stereotypes as well as varying degrees of
racial, religious, and age-based stereotypes in the corpus. A temporal analysis about word embeddings ~\cite{garg2018word} captures changes in gender and ethnic stereotypes over time.
Researchers attributed such problem partly to the biases in the datasets~\cite{jieyu,yao2017beyond} and word embeddings~\cite{garg2017word,caliskan2017semantics} but did not provide constructive solutions.

% \newcite{jieyu} observed that the gender bias problem was present in datasets for tasks like multi-label classification and visual semantic role labeling. 
% The gender bias would be further amplified in the models trained on these datasets. \newcite{yao2017beyond} showed that gender bias problems also existed in collaborative systems. They evaluated the bias in the system with several metrics.

%For example, \textit{GoogleNews} and \textit{Wikipedia} datasets are usually used to train word embeddings. Although these datasets seem to be comparatively unbiased, ~\newcite{kaiwei} showed that word embeddings trained on them still inherited gender stereotypes. \newcite{garg2017word} showed that word embeddings captured the temporal dynamics of gender and ethnic stereotypes in the society. \newcite{caliskan2017semantics} presented that female names were more likely to be associated with family words than male names. However they do not propose constructive solutions.




% In this work, they proposed a post-processing approach to reduce the gender bias, however as we point out previously, 
% However, few works have explored this problem before. 
% Previous work from~\cite{kaiwei}
% makes a comprehensive analysis on the geometry of gender bias in word embeddings and 
% provide a post-processing approach to remove the gender associations in the gender neutral words. 
% However, 
% the performance of this model highly depended on the accuracy of manually selected gender neutral word set. In addition, this method failed to keep the gender feature of each embedding and was not capable of interpreting the word representations. 

% The recent achievements on adversarial models motivated us on solving the gender debiasing problem. \cite{GAN}~\cite{unsuper}~\cite{ad_fl}~\cite{control} all contributed to the feature representation learning task in computer vision fields by using the adversarial techniques. In our paper, the minimax game was initially applied to debiase the word embeddings, and strongly helped the word representation interpretation.



